%% ma: L10-B12-0057
%% lop: 10
%% chuong: 5
%% bai: 12
%% dang: 10.12.4
%% loai: DS
%% muc: [TH, NB, TH, TH]
%% dap_an: "ĐĐSĐ"
%% nguon: "chuyen-de-cac-so-dac-trung-cua-mau-so-lieu-khong-ghep-nhom-toan-10-chuong-trinh-moi.pdf (GV. Phan Nhật Linh) (P01), Bài 01 Phần II Câu 1"
%% kien_thuc: "Lớp 10 Bài 12 (số gần đúng, sai số tuyệt đối, độ chính xác, sai số tương đối, quy tròn số)"
%% trang_thai: chinh-thuc
%% ghi_chu: "Ý a) gốc 'không vượt quá 0,663%' không chặt (0,5/75,4 ≈ 0,6631%), đã đổi thành 0,67%. Tự duyệt theo yêu cầu giáo viên 22/10/2026."
\begin{ex}
Kết quả đo chiều dài của một thửa đất là $75{,}4\text{ m}\pm0{,}5\text{ m}$ và đo chiều dài của một cây cầu là $466{,}2\text{ m}\pm0{,}5\text{ m}$. Xét tính đúng sai của các khẳng định sau:
\choiceTF{\True Đối với phép đo thửa đất, sai số tương đối không vượt quá $0{,}67\%$}{\True Đối với phép đo thửa đất, $\dfrac{d}{|a|}=\dfrac{0{,}5}{75{,}4}=\dfrac{5}{754}$}{Đối với phép đo cây cầu, sai số tương đối có thể lớn hơn $\dfrac{5}{4662}\approx0{,}107\%$}{\True Phép đo cây cầu có độ chính xác cao hơn phép đo thửa đất}
\loigiai{a) Đúng: $\delta\le\dfrac{0{,}5}{75{,}4}\approx0{,}663\%<0{,}67\%$.

b) Đúng: $\dfrac{0{,}5}{75{,}4}=\dfrac5{754}$.

c) Sai: $\delta\le\dfrac{0{,}5}{466{,}2}=\dfrac5{4662}\approx0{,}107\%$.

d) Đúng: $\dfrac5{4662}<\dfrac5{754}$ nên phép đo cây cầu chính xác hơn.}
\end{ex}
